% Abhakseite „Das kann ich“ – Zone nur im Gesamt (\mitzone)
%% TODO Ich-kann-Sätze: Platzhalter wie in den Aufgabendateien
\begin{abhakseite}
\ifdefined\mitzone
\abhakgruppe{Kennst du schon}
\abhak{1}{Flächenformeln der ebenen Figuren (Dreieck, Trapez, Parallelogramm, Raute/Drachen, Kreis) kennen und anwenden}
\abhak{2}{Volumen- und Oberflächenformeln der Elementarkörper (Prisma, Pyramide, Zylinder, Kegel) kennen}
\abhak{3}{Verbindungsvektoren bilden, Beträge berechnen, Mittelpunkte bestimmen}
\abhak{4}{Skalarprodukt und „null heißt senkrecht“}
\abhak{5}{Satz des Pythagoras für Seitenhöhen und fehlende Katheten, Strahlensätze und Ähnlichkeit für Streckfaktoren}
\abhak{6}{Gleichungen aus Formeln lösen (nach einer Größe umstellen, Wurzelgleichungen, quadratische Terme, Prozentrechnung für Anteile)}
\abhak{7}{Flächenformeln der ebenen Figuren (Dreieck, Trapez, Parallelogramm, Raute/Drachen, Kreis) kennen und anwenden – Fehler finden}
\abhak{8}{Flächenformeln der ebenen Figuren (Dreieck, Trapez, Parallelogramm, Raute/Drachen, Kreis) kennen und anwenden – selbst rechnen}
\fi
\abhakgruppe{Einheit 1 · Dreiecksflächen}
\abhak{9}{Welche Figur, welche Formel?}
\abhak{10}{Dreiecksflächen}
\abhak{11}{Dreiecksflächen – Fehler finden, Begründen, Darstellungswechsel, Anwendung}
\abhakgruppe{Einheit 2 · Vierecksflächen}
\abhak{12}{Vierecksflächen}
\abhak{13}{Ebene Figur: Flächeninhalt des Schattens eines geneigten Rechtecks bei senkrechtem Lichteinfall über eine Querschnittsskizze begründen}
\abhak{14}{Vierecksflächen – Fehler finden, Begründen, Darstellungswechsel, Anwendung}
\abhakgruppe{Einheit 3 · Einheit 3}
\abhak{15}{Volumen und Oberfläche}
\abhak{16}{Volumen und Oberfläche – Fehler finden, Begründen, Anwendung}
\abhakgruppe{Einheit 4 · Einheit 4}
\abhak{17}{Zusammengesetzt und Verhältnisse}
\abhak{18}{Zusammengesetzt und Verhältnisse – Fehler finden, Begründen, Darstellungswechsel, Anwendung}
\end{abhakseite}
